
% ============================================================
%  Buzzer Sub-Circuit: RPi GPIO → NPN transistor → Buzzer
%  Components: R1 (1kΩ), MMBT3904 (NPN), CEM-1203(42) buzzer,
%              1N4148W flyback diode
%
%  HOW TO COMPILE:
%    pdflatex buzzer_circuit.tex
%  Requires: circuitikz package (part of TeX Live / MiKTeX)
% ============================================================

\documentclass[border=10pt]{standalone}
% 'standalone' produces a tight bounding box around the drawing.
% 'border=10pt' adds a small margin around the content.

\usepackage{circuitikz}
% The main package for drawing electronic circuits in LaTeX.
% It builds on TikZ and provides component shapes (resistors,
% transistors, diodes, etc.) via \draw ... to[component] ...

\usepackage{xcolor}
% Lets us name colors for annotation text.

% ── Convenience color for annotation comments ──────────────
\definecolor{annotate}{RGB}{0,100,180}

\begin{document}
	

% ----------------------------------------------------------
% \begin{circuitikz}[...]
%   Opens the drawing environment.
%   american  → uses American resistor symbol (zigzag)
%   european  → would use IEC rectangle style
%   scale / transform shape scale → resize all components
% ----------------------------------------------------------
\begin{circuitikz}[european, scale=1.1, transform shape]

	% Coordinates grid
	\def\xIN{0}       % left edge / GPIO pin
	\def\xNPN{4}        % NPN centre column
	\def\xBuzz{4}       % buzzer / diode column
	\def\xDiode{2}    % flyback diode column
	\def\xR{3}
	
	\def\yVBUS{6}        % main horizontal bus  (R1 / NPN base)
	\def\yBuzztop{5.5}        % +5V rail
	\def\yBuzzbot{3}        % +5V rail
	\def\yGPIO{1.5}        % +5V rail
	\def\yGND{0}       % ground level

	% ── GND Rail ────────────────────────────────────────────────────
	\draw (\xIN, \yGND)
	node[left]{\textbf{GND}};
	\node[circ] at (\xNPN, \yGND) {};
	
	% ── GND rail: horizontal thick line ─────────────────────
	\draw[very thick] (\xIN, \yGND) to[short, o-*] (\xBuzz, \yGND);
	
	% ── npn transistor named T1 ─────────────────────
	\node[npn, xscale=1] (T1) at (\xBuzz, \yGPIO) {};

	% Label the transistor
	\node[right=4pt] at (T1.center) {MMBT3904};
	
	% ── GPIO to Resistor to T1.base─────────────────────
	\draw (\xIN, \yGPIO)
	node[left]{\textbf{Rpi GPIO (PWM)}}
	to[R, l=R1, a=1\,k$\Omega$, o-] (T1.B);

	
	% ── 5. EMITTER → GND ────────────────────────────────────
	\draw (T1.E) to[short] (\xNPN, \yGND);  % vertical drop

	% ── 5. COLLECTOR → Buzzer ───────────────────────────────
    \draw (T1.C) to[short, -*] (\xNPN, \yBuzzbot);  % vertical drop

	% ── 5. Buzzer ───────────────────────────────
	\draw (\xNPN, \yBuzzbot) to 
	[buzzer,                        % buzzer component
	l_={\parbox{3cm}{\centering Piezo Buzzer \\ CEM-1203(42)}}]
	(\xNPN, \yBuzztop);
	% Wire from resistor junction to transistor Base
	% T1.B  → the Base terminal anchor of node T1
	%\draw (3,0) -- (T1.B);



\end{circuitikz}
\end{document}	
	
	% ── GPIO, R1, npn ───────────────────────────────────────────────
% ── Junction dot at NPN emitter drop-down point ─────────
%\node[circ] at (\xNPN, \yGND) {};
	 	
	

	% ── 5V Rail ─────────────────────────────────────────────────────
	\draw (\xIN, \yVBUS)
	node[left]{\textbf{+5V}};	
	\draw[very thick] (\xIN, \yGND) -- (\xDiode, \yVBUS) o-*;
	\draw[very thick] (\xDiode, \yVBUS) -- (\xBuzz, \yVBUS)  -*;

\end{circuitikz}
\end{document}	
	% ── COORDINATE LAYOUT (all distances in cm) ─────────────
	% We place named coordinates first so that wires and
	% components can reference them by name.
	%
	%  RPi ──R1──┬── Base(B)
	%            │
	%       Emitter(E) ── GND
	%  +5V ─ Buzzer ──┬── Collector(C)
	%                 │
	%            1N4148W (flyback, anode→C, cathode→+5V)
	
	% ── 1. INPUT LABEL: RPi GPIO (PWM) ──────────────────────
	% \node places text or a shape at a coordinate.
	% [left] aligns the anchor to the left of the coordinate.
	\node[left] at (0,0) {\textbf{RPi GPIO (PWM)}};
	
	% ── 2. SERIES RESISTOR R1 ───────────────────────────────
	% \draw begins a path.
	% (0,0) is the start coordinate.
	% -- is a straight wire segment to the next point or component.
	% to[R, ...] draws a resistor along the path.
	%   R        → resistor component type
	%   l=R1     → label above the component (l=label)
	%   a=1\,k$\Omega$ → label below (a=annotation)
	%   *-       → draw a filled dot at the start terminal
	%   -*       → filled dot at the end terminal
	%   (use - with no * for no dot)
	\draw (0,0)
	to[R, l=R1, a=1\,k$\Omega$, o-] (3,0);
	%       ↑ 'o-' draws an open circle at the left terminal
	%         (useful for external connection points / pins)
	
	% ── 3. JUNCTION NODE at base-resistor join ───────────────
	% After the resistor we mark the junction explicitly.
	% 'node[circ]' draws a filled dot (junction marker).
	\draw (3,0) node[circ]{};
	
	% ── 4. NPN TRANSISTOR (MMBT3904) ─────────────────────────
	% 'npn' is the CircuiTikZ keyword for an NPN BJT.
	% The transistor is placed with its CENTER at (4,0) but
	% CircuiTikZ anchors it by the *base* pin by default when
	% using the 'at (x,y)' syntax inside a node.
	%
	% Named node syntax:  \node[npn] (T1) at (x,y) {};
	%   (T1) → gives the transistor the name "T1" so we can
	%           reference its terminals: T1.B, T1.C, T1.E
	\node[npn, xscale=1] (T1) at (4.5, 0) {};
	% xscale=1 → no horizontal flip; use xscale=-1 to mirror
	
	% Label the transistor
	\node[right=4pt] at (T1.center) {MMBT3904};
	
	% Wire from resistor junction to transistor Base
	% T1.B  → the Base terminal anchor of node T1
	\draw (3,0) -- (T1.B);
	
	% ── 5. EMITTER → GND ────────────────────────────────────
	% T1.E → Emitter terminal of T1
	% to[short, -*] draws a plain wire with a filled dot at end
	% 'node[ground]{}' places the ground symbol (GND)
	\draw (T1.E)
	to[short] (T1.E |- 0,-2.5)  % vertical drop
	node[ground]{};             % GND symbol at the bottom
	% 'node[ground]' draws the three-line GND symbol.
	% 'T1.E |- 0,-3.5' means: same X as T1.E, same Y as (0,-3.5)
	% — a TikZ coordinate shorthand for vertical/horizontal intersections.

	% ── 6. COLLECTOR WIRE up to buzzer ───────────────────────
	% T1.C → Collector terminal of T1
	% We'll run the collector wire upward to (4.5, 0.5).
 	\coordinate (COL) at (4.5, 0.5);  % named coord for reuse
 	\draw (T1.C) -- (COL);
	
	% ── 7. FLYBACK DIODE 1N4148W (across buzzer, in parallel) ─
	% 'D' → generic diode; CircuiTikZ draws it with the standard
	%       triangle-and-bar symbol.
	% Direction matters: current flows from anode (+) to cathode (−)
	% in forward bias.  Flyback diode is reverse-biased normally;
	% anode at Collector side, cathode at +5V side.
	%
	% to[D, invert] → 'invert' flips the diode direction so the
	%                 arrow points FROM cathode TO anode in the
	%                 drawing path direction.  Here we draw the
	%                 path upward (COL → 5V rail) so 'invert'
	%                 puts the anode at the bottom (Collector).
	\draw (COL)
	to[D, invert, l=1N4148W, color=red] (4.5, 2.5);
	%       ↑ 'color=red' colors just this component for visibility
	
	% ── 8. +5V RAIL and BUZZER ──────────────────────────────
	% We draw the +5V rail from the left, through the buzzer,
	% down to the Collector/diode node.
	%
	% 'buzzer' is a built-in CircuiTikZ component type.
	% It draws the IEC piezo-buzzer symbol.
	\draw (0, 2.5)
	node[left]{\textbf{+5V}}          % +5V label at left
	to[short, o-] (1.5, 2.5)          % short wire from pin
	to[buzzer,                         % buzzer component
	l=CEM-1203(42),                 % label above
	a=Piezo Buzzer]                 % annotation below
	(4.5, 2.5);                     % end at diode/collector junction
	
	% Filled junction dot where buzzer meets diode/collector rail
	\draw (4.5, 2.5) node[circ]{};
	
	% Vertical wire connecting +5V end down to COL junction
	% (the diode already covers this path; just confirm node)
	% Already connected through the diode draw above.
	
	% ── 9. ANNOTATION LABELS for terminals ──────────────────
	% \node with text only, placed near key points.
	\node[annotate, above=2pt] at (T1.B) {\footnotesize Base (B)};
	\node[annotate, right=6pt] at (T1.C) {\footnotesize Collector (C)};
	\node[annotate, right=6pt] at (T1.E) {\footnotesize Emitter (E)};
	
	% ── 10. GROUND LABEL ─────────────────────────────────────
	\node[below=2pt] at (T1.E |- 0,-3.5) {\footnotesize GND};
	
\end{circuitikz}
% ----------------------------------------------------------
% END OF DRAWING
%
% KEY CircuiTikZ CONCEPTS USED:
%  \node[component](name) at (x,y){}  → place a component node
%  \draw (a) to[component, options] (b) → draw component on wire
%  to[R]    → resistor
%  to[D]    → diode  (add 'invert' to flip direction)
%  to[buzzer] → piezo buzzer symbol
%  node[npn]  → NPN BJT transistor
%  node[ground] → ground symbol
%  node[circ]   → filled junction dot
%  l=...        → label (above/left of component)
%  a=...        → annotation (below/right of component)
%  o-  / -o     → open terminal circle (external pin)
%  *-  / -*     → filled dot at terminal
%  T1.B / T1.C / T1.E → named anchor points on transistor T1
%  (A |- B)     → TikZ: x of A, y of B (right-angle helper)
% ----------------------------------------------------------
\end{figure}
\end{document}